\markdownRendererDocumentBegin
\markdownRendererWarning{The `hybrid` option has been soft-deprecated.}{The `hybrid` option has been soft-deprecated.}{Consider using one of the following better options for mixing TeX and markdown: `contentBlocks`, `rawAttribute`, `texComments`, `texMathDollars`, `texMathSingleBackslash`, and `texMathDoubleBackslash`. For more information, see the user manual at <https://witiko.github.io/markdown/>.}{Consider using one of the following better options for mixing TeX and markdown: `contentBlocks`, `rawAttribute`, `texComments`, `texMathDollars`, `texMathSingleBackslash`, and `texMathDoubleBackslash`. For more information, see the user manual at <https://witiko.github.io/markdown/>.}\markdownRendererSectionBegin
\markdownRendererHeadingOne{Guía Rápida de Comandos (Cheat Sheet)}\markdownRendererInterblockSeparator
{}Esta guía contiene los comandos esenciales para gestionar la infraestructura, la conexión y el entrenamiento del SLM.\markdownRendererInterblockSeparator
{}\markdownRendererSectionBegin
\markdownRendererHeadingTwo{1. Gestión de Infraestructura (Terraform)}\markdownRendererInterblockSeparator
{}\markdownRendererEmphasis{Ejecutar desde la carpeta \markdownRendererCodeSpan{app/terraform/} en tu ordenador local.}\markdownRendererInterblockSeparator
{}\markdownRendererUlBeginTight
\markdownRendererUlItem \markdownRendererStrongEmphasis{Inicializar}: \markdownRendererCodeSpan{terraform init}\markdownRendererUlItemEnd 
\markdownRendererUlItem \markdownRendererStrongEmphasis{Ver cambios pendientes}: \markdownRendererCodeSpan{terraform plan}\markdownRendererUlItemEnd 
\markdownRendererUlItem \markdownRendererStrongEmphasis{Desplegar cambios}: \markdownRendererCodeSpan{terraform apply -auto-approve}\markdownRendererUlItemEnd 
\markdownRendererUlItem \markdownRendererStrongEmphasis{Obtener la IP de la máquina}: \markdownRendererCodeSpan{terraform output ec2\markdownRendererUnderscore{}public\markdownRendererUnderscore{}ip}\markdownRendererUlItemEnd 
\markdownRendererUlItem \markdownRendererStrongEmphasis{Destruir todo (Ahorro de costes)}: \markdownRendererCodeSpan{terraform destroy -auto-approve}\markdownRendererUlItemEnd 
\markdownRendererUlItem \markdownRendererStrongEmphasis{Importar recurso existente}: \markdownRendererCodeSpan{terraform import <direccion\markdownRendererUnderscore{}recurso> <id\markdownRendererUnderscore{}en\markdownRendererUnderscore{}aws>}\markdownRendererUlItemEnd 
\markdownRendererUlEndTight \markdownRendererInterblockSeparator
{}
\markdownRendererSectionEnd \markdownRendererSectionBegin
\markdownRendererHeadingTwo{2. Conexión Remota (SSH)}\markdownRendererInterblockSeparator
{}\markdownRendererEmphasis{Ejecutar desde la carpeta donde tengas el archivo \markdownRendererCodeSpan{.pem}.}\markdownRendererInterblockSeparator
{}\markdownRendererUlBeginTight
\markdownRendererUlItem \markdownRendererStrongEmphasis{Configurar permisos (solo la primera vez)}: \markdownRendererCodeSpan{chmod 400 tfm-slm.pem}\markdownRendererUlItemEnd 
\markdownRendererUlItem \markdownRendererStrongEmphasis{Conectar a la instancia}: \markdownRendererCodeSpan{ssh -i "tfm-slm.pem" ec2-user@<IP\markdownRendererUnderscore{}PUBLICA>}\markdownRendererUlItemEnd 
\markdownRendererUlItem \markdownRendererStrongEmphasis{Limpiar IP antigua (si da error de host)}: \markdownRendererCodeSpan{ssh-keygen -R <IP\markdownRendererUnderscore{}PUBLICA>}\markdownRendererUlItemEnd 
\markdownRendererUlEndTight \markdownRendererInterblockSeparator
{}
\markdownRendererSectionEnd \markdownRendererSectionBegin
\markdownRendererHeadingTwo{3. Control del Entrenamiento (Docker)}\markdownRendererInterblockSeparator
{}\markdownRendererEmphasis{Ejecutar una vez dentro de la instancia EC2 por SSH.}\markdownRendererInterblockSeparator
{}\markdownRendererUlBeginTight
\markdownRendererUlItem \markdownRendererStrongEmphasis{Ver contenedores activos}: \markdownRendererCodeSpan{docker ps}\markdownRendererUlItemEnd 
\markdownRendererUlItem \markdownRendererStrongEmphasis{Ver logs del entrenamiento (en vivo)}: \markdownRendererCodeSpan{docker logs -f <ID\markdownRendererUnderscore{}CONTENEDOR>}\markdownRendererUlItemEnd 
\markdownRendererUlItem \markdownRendererStrongEmphasis{Ver todos los contenedores (incluidos fallidos)}: \markdownRendererCodeSpan{docker ps -a}\markdownRendererUlItemEnd 
\markdownRendererUlItem \markdownRendererStrongEmphasis{Ver logs de un contenedor que ya falló}: \markdownRendererCodeSpan{docker logs <ID\markdownRendererUnderscore{}CONTENEDOR>}\markdownRendererUlItemEnd 
\markdownRendererUlItem \markdownRendererStrongEmphasis{Borrar todos los contenedores parados}: \markdownRendererCodeSpan{docker container prune -f}\markdownRendererUlItemEnd 
\markdownRendererUlEndTight \markdownRendererInterblockSeparator
{}
\markdownRendererSectionEnd \markdownRendererSectionBegin
\markdownRendererHeadingTwo{4. Monitorización de Hardware (NVIDIA)}\markdownRendererInterblockSeparator
{}\markdownRendererEmphasis{Ejecutar una vez dentro de la instancia EC2 por SSH.}\markdownRendererInterblockSeparator
{}\markdownRendererUlBeginTight
\markdownRendererUlItem \markdownRendererStrongEmphasis{Ver estado de la GPU y VRAM}: \markdownRendererCodeSpan{nvidia-smi}\markdownRendererUlItemEnd 
\markdownRendererUlItem \markdownRendererStrongEmphasis{Monitorizar GPU cada segundo}: \markdownRendererCodeSpan{nvidia-smi -l 1}\markdownRendererUlItemEnd 
\markdownRendererUlEndTight \markdownRendererInterblockSeparator
{}
\markdownRendererSectionEnd \markdownRendererSectionBegin
\markdownRendererHeadingTwo{5. Depuración de Arranque (Cloud-Init)}\markdownRendererInterblockSeparator
{}\markdownRendererEmphasis{Si la máquina acaba de encenderse y \markdownRendererCodeSpan{docker ps} no muestra nada.}\markdownRendererInterblockSeparator
{}\markdownRendererUlBeginTight
\markdownRendererUlItem \markdownRendererStrongEmphasis{Ver progreso de instalación inicial}: \markdownRendererCodeSpan{sudo tail -f /var/log/cloud-init-output.log}\markdownRendererUlItemEnd 
\markdownRendererUlEndTight \markdownRendererInterblockSeparator
{}
\markdownRendererSectionEnd \markdownRendererSectionBegin
\markdownRendererHeadingTwo{6. Gestión de Sesiones (Pausar y Reanudar)}\markdownRendererInterblockSeparator
{}\markdownRendererEmphasis{Permite ahorrar costes deteniendo la GPU sin perder el progreso del entrenamiento.}\markdownRendererInterblockSeparator
{}\markdownRendererUlBegin
\markdownRendererUlItem \markdownRendererStrongEmphasis{Pausar Entrenamiento (Manteniendo Checkpoints en S3)}:\markdownRendererSoftLineBreak
{}\markdownRendererCodeSpan{terraform destroy -target=aws\markdownRendererUnderscore{}spot\markdownRendererUnderscore{}instance\markdownRendererUnderscore{}request.training -auto-approve}\markdownRendererSoftLineBreak
{}\markdownRendererEmphasis{(Esto destruye la máquina pero mantiene el bucket de S3 intacto).}\markdownRendererUlItemEnd 
\markdownRendererUlItem \markdownRendererStrongEmphasis{Reanudar Entrenamiento}:\markdownRendererSoftLineBreak
{}\markdownRendererCodeSpan{terraform apply -auto-approve}\markdownRendererSoftLineBreak
{}\markdownRendererEmphasis{(La nueva máquina detectará automáticamente el checkpoint en S3 y continuará desde la época actual).}\markdownRendererUlItemEnd 
\markdownRendererUlEnd 
\markdownRendererSectionEnd 
\markdownRendererSectionEnd \markdownRendererDocumentEnd